
\chapter[This is Chapter One Title containing authors and affiliations]{This is Chapter One Title containing authors and affiliations\protect\footnote{This is a sample for chapter footnote.}}

\author*[1]{Author I}
\author[2]{Author II}
\author[1]{Author III}
\author[3]{Author IV}

\address[1]{\orgdiv{I Author Organization Division Name}, 
\orgname{Organization Name}, 
\postcode{Postal Code}, \countrypart{Part of the Country}, 
     \city{City Name}, \street{Street Name}, \country{Country}}%

\address[2]{\orgdiv{II Author Organization Division Name}, 
\orgname{Organization Name}, 
\postcode{Postal Code}, \countrypart{Part of the Country}, 
     \city{City Name}, \street{Street Name}, \country{Country}}%

\address[3]{\orgdiv{II Author Organization Division Name}, 
\orgname{Organization Name}, 
\postcode{Postal Code}, \countrypart{Part of the Country}, 
     \city{City Name}, \street{Street Name}, \country{Country}}%

\address*{Corresponding Author: Author; \email{corresauthor@gmail.com}}

\maketitle% This tag is required to print author and address in the output

\begin{abstract}{Abstract Title Here}
Abstract para text. Abstract para text. Abstract para text. Abstract para text. 
Abstract para text. Abstract para text. Abstract para text. Abstract para text.

Abstract next para text. Abstract para text. Abstract para text. Abstract para text. 
Abstract para text. If multiple abstracts exists, then \textbf{absgroup} tag 
will be inserted by the tool.
\end{abstract}

\begin{abstract}{}
Abstract text. Abstract text. Abstract text. Abstract text. 
\end{abstract}

\keywords{Key word 1, Key word 2, \ldots}

\section{This is First Level Heading}

\lipsum[1-2]

\subsection{This is Second Level Heading}

\lipsum[3]

For example, multiple citations from the \index{bibliography} of this \index{article}:
citet: \citet{CR7,CR8}, citep: \citep{CR9,CR6}.
As you see in Table~\ref{tab1-2}, the citations are \index{automatically!hyperlinked!citations} to their
reference in the bibliography. Refer Figure~\ref{fig1-1} for details:

\begin{figure}
\includegraphics{01.eps}
\caption{Figure Caption. Figure Caption.
Figure Caption. Figure Caption. Figure Caption. 
Figure Caption. Figure Caption.
\label{fig1-1}}
\end{figure}

For example, multiple citations from the \index{bibliography} of this \index{article}:
citet: \citet{CR9,CR6}, citep: \citep{CR9,CR6}.
As you see in Table~\ref{tab1-1}, the citations are \index{main entry with turn over lines main entry with turn over lines!subentry with turn over lines subentry with turn over lines!subsubentry with turn over lines subsubentry with turn over lines} to their
reference in the bibliography (Equation~\ref{eq1-1}).
\begin{equation}
\mathcal{L}\quad \mathbf{\mathcal{L}} = i \bar{\psi} \gamma^\mu D_\mu \psi - 
\frac{1}{4} F_{\mu\nu}^a F^{a\mu\nu} - m \bar{\psi} \psi\label{eq1-1}
\end{equation}

\subsubsection{This is Third Level Heading}

The manifestation of solar activity (flares, bursts, and others) occurs over the whole Sun, and most of radio astronomy observations are made from the Earth's surface, whereas a significant part of solar radio events (those from the far side of the Sun) is not available for terrestrial observers.
\begin{equation}
\mathcal{L}\quad \mathbf{\mathcal{L}} = i \bar{\psi} \gamma^\mu D_\mu \psi - 
\frac{1}{4} F_{\mu\nu}^a F^{a\mu\nu} - m \bar{\psi} \psi\label{eq1-2}
\end{equation}

\paragraph{This is Fourth Level Heading}

\lipsum[5]

\begin{table}
\caption{Enter table caption here.\label{tab1-1}}{%
\begin{tabular}{@{}cccc@{}}
\toprule
Tap     &Relative   &Relative   &Relative mean\\
number  &power (dB) &delay (ns) &power (dB)\\
\midrule
3 &0$-9.0$  &68,900\footnotemark[1] &$-12.8$\\
4 &$-10.0$ &12,900\footnotemark[2] &$-10.0$\\
5 &$-15.0$ &17,100 &$-25.2$\\
\botrule
\end{tabular}}{\footnotetext[]{Source: Example for table source text.}
\footnotetext[1]{Example for a first table footnote. Example for a first table footnote. Example for a first table footnote. Example for a first table footnote.}
\footnotetext[2]{Example for a second table footnote.}}
\end{table}

\subparagraph{This is Fifth Level Heading}

\lipsum[6]

The manifestation of solar activity (flares, bursts, and others) occurs over the whole Sun, and most of radio astronomy observations are made from the Earth's surface, whereas a significant part of solar radio events (those from the far side of the Sun) is not available for terrestrial observers (Equation~\ref{eq1-2}).

\begin{table}
\caption{Enter table caption here.\label{tab1-2}}{%
\begin{tabular}{@{}cccc@{}}
\toprule
Tap     &Relative   &Relative   &Relative mean\\
number  &power (dB) &delay (ns) &power (dB)\\
\midrule
3 &0$-9.0$  &68,900\footnotemark[1] &$-12.8$\\
4 &$-10.0$ &12,900\footnotemark[2] &$-10.0$\\
5 &$-15.0$ &17,100 &$-25.2$\\
\botrule
\end{tabular}}{\footnotetext[]{Source: Example for table source text.}
\footnotetext[1]{Example for a first table footnote. Example for a first table footnote. Example for a first table footnote. Example for a first table footnote.}
\footnotetext[2]{Example for a second table footnote.}}
\end{table}

\backmatter

\bibliographystyle{plainnat}
\bibliography{wiley}%

